\section{Schlussbemerkungen}
\label{sec:schluss}

Eine fehlende Zeile in einer Basisklasse, auf den ersten Blick ein
Entwicklungsfehler wie viele andere. Tatsächlich war die fehlende
Mandantenfilterung in \texttt{BaseService.get()} kein isoliertes Problem an
einem einzelnen Endpunkt, sondern ein systemischer Konstruktionsfehler, der
sich auf nahezu alle CRUD-Endpunkte mit organisationsgebundenen Ressourcen
auswirkte. Und weil in derselben Methode auch der \texttt{deleted\_at}-Filter
fehlte, liessen sich soft-gelöschte Datensätze per direkter UUID-Abfrage
weiterhin lesen. Ein Soft-Delete-Bypass, der ohne diese Analyse vermutlich
unentdeckt geblieben wäre.

Die Behebung erfolgte bewusst an der gemeinsamen Wurzel des Problems. Statt
einzelne Endpunkte punktuell nachzubessern, wurde die Basisklasse um einen
optionalen \texttt{organization\_id}-Parameter sowie einen konsistenten
Soft-Delete-Filter erweitert und anschliessend an allen relevanten Aufrufern
korrekt verdrahtet. Dieses Vorgehen reduzierte den Änderungsaufwand pro
Endpunkt, verbesserte die Konsistenz der Sicherheitslogik und senkte das Risiko,
dass künftig neue Router erneut ohne Mandantenscoping implementiert werden.

Das Testergebnis spricht für sich. Die zuvor als
Proof-of-Concept formulierten Angriffe schlagen nach der Anpassung fehl, die neu
hinzugefügten Regressionstests decken beide Schwachstellenklassen gezielt ab,
und die bestehende Testsuite läuft nach Anpassung der betroffenen Alt-Tests
wieder vollständig durch. Der finale Stand mit \texttt{187 passed, 2 xfailed}
zeigt, dass der Sicherheitsfix nicht nur lokal funktioniert, sondern auch mit
dem restlichen Backend vereinbar ist.

Multi-Tenant-Isolation ist kein Komfortmerkmal, sie ist eine
Systemeigenschaft. Sobald Ressourcen per UUID adressierbar sind, muss die
Autorisierung bei jedem Zugriff explizit an den aktuellen Mandanten
gebunden werden. Wird diese Prüfung an einer zentralen Abstraktion
vergessen, skaliert die Schwachstelle mit der Wiederverwendung genauso
stark wie der Architekturvorteil selbst. Insofern ist das beobachtete
Muster eine klassische Verletzung des von Saltzer und Schroeder
formulierten Prinzips \enquote{Complete Mediation}, das verlangt, jeden
Zugriff auf ein geschütztes Objekt einzeln zu autorisieren
\parencite{saltzer_schroeder_1975}.

